\bookchapter{Capitolo 30: Casting, presentazioni e connettività auto}

\chapterintro{Lo schermo del telefono, per quanto grande, arriva a 7 pollici; ma la tua TV ne ha 55 e il proiettore 100. “Gettare” lo schermo del telefono su uno grande cambia dimensione all'esperienza. Questo capitolo illustra tre scenari ad alta frequenza — casting, presentazioni e connettività auto — per aiutarti a estendere le capacità del telefono a schermi più grandi.}

\booksection{Le tre famiglie del casting wireless — non confonderle}

Il casting wireless sembra sempre “schermo del telefono alla TV”, ma le tecnologie sottostanti sono tre e non interoperabili:

\begin{center}
\begin{tabularx}{\linewidth}{|X|X|X|}
\hline
\textbf{Tecnologia} & \textbf{Chi la usa} & \textbf{Caratteristiche} \tabularnewline
\hline
AirPlay & iPhone/iPad/Mac & Ecosistema Apple, l'esperienza più fluida. Richiede Apple TV o una TV che supporti AirPlay \tabularnewline
Miracast & Telefoni Android & Wi-Fi Direct diretto, non serve router. Il "casting wireless" di Android di default usa questo \tabularnewline
DLNA & Tutte le piattaforme & Trasmette solo file multimediali (video/musica/foto), non replica lo schermo. Ottimo per inviare video alla TV \tabularnewline
\hline
\end{tabularx}
\end{center}

\booksubsection{Metodi operativi per piattaforma:}

\begin{itemize}
  \item \textbf{iPhone}: Centro di Controllo > Screen Mirroring > scegli la tua TV/Apple TV. Se la TV non supporta AirPlay, puoi comprare un set-top box Apple TV o una chiavetta TV che supporti AirPlay.
  \item \textbf{Android}: abbassando il pannello notifiche > trova il pulsante "Casting wireless"/"Casting intelligente"/"Multi-screen interattivo" > scegli la TV. Telefono e TV devono essere sulla stessa rete Wi-Fi.
  \item \textbf{Soluzione universale}: se la TV è troppo vecchia e non supporta alcun casting wireless, compra un Chromecast o una Xiaomi Box/Huawei Box da collegare alla porta HDMI della TV, e potrai ricevere il casting.
\end{itemize}

\booksection{Casting via cavo — più stabile, latenza zero}

Il casting wireless è comodo ma ha latenza, poco adatto a giochi e presentazioni. Il casting via cavo risolve:

\begin{itemize}
  \item \textbf{iPhone}: serve un cavo convertitore Lightning/USB-C to HDMI (circa 100-200 yuan). Collegato, lo schermo del telefono esce direttamente sulla porta HDMI della TV/proiettore.
  \item \textbf{Android}: alcuni telefoni supportano USB-C to HDMI diretto (serve il supporto a DP Alt Mode, presente nella maggior parte dei top di gamma). Basta un cavo USB-C to HDMI.
  \item \textbf{Scenario sala riunioni}: se fai spesso presentazioni con il telefono, tieni in borsa un cavo convertitore HDMI. Le TV/proiettori delle sale riunioni hanno quasi tutti una porta HDMI: basta collegare e usarle, molto più affidabile del casting wireless.
\end{itemize}

\booksection{Il telefono diventa computer — modalità desktop}

Alcuni telefoni Android supportano la “modalità desktop” — collegandoli a un monitor, il telefono emette un'interfaccia simile a un desktop, utilizzabile con mouse e tastiera:

\begin{itemize}
  \item \textbf{Samsung DeX}: la funzione di punta dei top di gamma Samsung. Collegato il monitor, compare un'interfaccia desktop, con multi-finestra; l'esperienza più vicina a un vero computer. Con tastiera e mouse Bluetooth, in viaggio puoi non portare il portatile.
  \item \textbf{Modalità desktop Huawei}: i Huawei Mate/Pura, collegati al monitor, hanno una funzione simile. Supporta multi-finestra, file manager, App di produttività.
  \item \textbf{Motorola Ready For}: anche i top di gamma Moto la supportano, simile a DeX.
  \item \textbf{Scenari adatti}: leggero lavoro d'ufficio (email, modifiche documenti), presentazioni PPT, video su grande schermo. Non adatto a uso intensivo di Excel o progettazione professionale.
\end{itemize}

\booksection{Connettività auto — il modo corretto di usare il telefono in macchina}

Usare il telefono alla guida è illegale, ma navigazione, musica e chiamate sono necessità. I sistemi di connettività auto consentono di usare le funzioni chiave del telefono sullo schermo dell'auto, mentre il telefono resta da parte:

\begin{itemize}
  \item \textbf{Apple CarPlay}: la prima scelta per gli utenti iPhone. Collegando via USB o wireless, sullo schermo dell'auto compare un'interfaccia semplificata dell'iPhone — navigazione (Amap/Baidu Maps), musica, telefono, messaggi, podcast. Il controllo vocale Siri libera le mani.
  \item \textbf{Android Auto}: la soluzione per Android. Funzioni simili a CarPlay, supporta Google Maps, Spotify ecc. In Cina, a causa delle restrizioni dei servizi Google, alcune funzioni non sono disponibili.
  \item \textbf{Huawei HiCar}: soluzione proprietaria Huawei, alta percentuale di supporto su modelli cinesi. Oltre a navigazione e musica, offre funzioni speciali come videocall Huawei e integrazione telecamera di parcheggio.
  \item \textbf{Baidu CarLife}: soluzione cross-platform, disponibile sia su iOS sia su Android. Meno funzioni ma alta compatibilità: molte auto vecchie supportano solo CarLife.
\end{itemize}

\begin{tcbwarn}{Promemoria di sicurezza}
Qualsiasi soluzione di connettività auto, alla guida usa per quanto possibile il controllo vocale. Imposta la destinazione di navigazione in anticipo, seleziona la playlist prima di partire, non operare lo schermo dell'auto in marcia.

\end{tcbwarn}

\booksection{Casting in sala riunioni — niente più confusione}

Le riunioni sono lo scenario più frequente di casting, e anche il più problematico. Ecco alcune esperienze pratiche:

\begin{itemize}
  \item \textbf{Arriva in anticipo per testare}: non aspettare l'inizio della riunione per collegarti. Arriva 10 minuti prima e verifica la modalità di collegamento. Il casting Wi-Fi fallisce più spesso di quanto pensi.
  \item \textbf{Tieni pronto un piano B via cavo}: in borsa un cavo convertitore HDMI. Quando il wireless non si collega, il cavo è la salvezza.
  \item \textbf{Disattiva l'anteprima delle notifiche}: durante il casting, i tuoi messaggi WeChat appaiono sul grande schermo della sala. Prima di proiettare, disattiva nelle impostazioni l'"anteprima contenuti notifiche schermata di blocco", o attiva la "modalità presentazione" (alcuni telefoni la supportano, nasconde le notifiche durante il casting).
  \item \textbf{Usa "documento presentato" non "telefono operato"}: se devi mostrare un PPT, salvalo sul telefono e aprilo a schermo intero con WPS, invece di proiettare l'intero desktop del telefono. Più professionale e più fluido.
\end{itemize}

\begin{tcbpractice}{Fai: completa un casting dal telefono alla TV}
\begin{enumerate}
  \item Assicurati che telefono e TV siano sulla stessa rete Wi-Fi.
  \item Apri sul telefono una foto o un video.
  \item Prova il casting wireless (iPhone con Screen Mirroring, Android con il pulsante Casting wireless).
  \item Se fallisce, verifica che la TV supporti il protocollo corrispondente, o prova DLNA (premi il pulsante "Casting" nell'App Galleria/Video).
  \item Dopo il successo, prova altre operazioni sul telefono per vedere se lo schermo grande si aggiorna in sincrono.
\end{enumerate}

\end{tcbpractice}

\begin{tcbsummary}{Riepilogo del capitolo}
\begin{itemize}
  \item Tre famiglie di casting wireless: AirPlay (Apple), Miracast (Android), DLNA (cross-platform per media)
  \item Il casting via cavo è più stabile e a latenza zero; per le sale riunioni è indispensabile un cavo convertitore HDMI
  \item Samsung DeX / modalità desktop Huawei trasformano il telefono in un computer leggero
  \item Connettività auto: quattro soluzioni — CarPlay/Android Auto/HiCar/CarLife
  \item Regole d'oro del casting in riunione: test in anticipo, piano B pronto, disattivare l'anteprima notifiche
\end{itemize}

\end{tcbsummary}
